\section{Granulometria, textura e morfologia dos grãos}

\begin{definicao}
\textbf{Granulometria} é a descrição da distribuição dos tamanhos das partículas de um sedimento ou rocha sedimentar. Ela pode ser determinada por peneiramento, sedimentação, difração a laser ou análise de imagens, conforme a faixa de tamanhos e o objetivo do estudo.
\end{definicao}

\subsection{Escala de Udden--Wentworth}

A escala de Udden--Wentworth organiza os sedimentos clásticos em classes cujos limites variam por fatores de dois \parencite{wentworth1922}. Os limites fundamentais são:

\begin{table}[H]
  \centering
  \caption{Classes granulométricas fundamentais.}
  \label{tab:wentworth-resumo}
  \begin{tabular}{lll}
    \toprule
    \textbf{Fração} & \textbf{Diâmetro aproximado} & \textbf{Rocha consolidada} \\
    \midrule
    Cascalho & maior que \SI{2}{\milli\metre} & conglomerado ou brecha \\
    Areia    & \SIrange{0.0625}{2}{\milli\metre} & arenito \\
    Silte    & \SIrange{0.0039}{0.0625}{\milli\metre} & siltito \\
    Argila   & menor que \SI{0.0039}{\milli\metre} & argilito \\
    \bottomrule
  \end{tabular}
\end{table}

Na escala phi, usada em sedimentologia, o tamanho é expresso por

\[
\phi=-\log_2 d,
\]

em que \(d\) é o diâmetro da partícula em milímetros. Assim, valores maiores de \(\phi\) representam partículas menores.

\subsection{Seleção granulométrica}

A \textbf{seleção} indica a dispersão dos tamanhos dos grãos. Um sedimento bem selecionado concentra-se em poucas classes granulométricas; um sedimento mal selecionado contém ampla variedade de tamanhos. A seleção depende do agente, da energia e da duração do transporte, mas também pode ser herdada da área-fonte.

Quando a distribuição é quantificada em unidades phi, podem ser empregadas classes como muito bem selecionada, bem selecionada, moderadamente selecionada, mal selecionada e muito mal selecionada. O método estatístico e os limites adotados devem ser declarados, pois classificações distintas coexistem na literatura.

\subsection{Arredondamento e esfericidade}

\textbf{Arredondamento} e \textbf{esfericidade} descrevem propriedades diferentes:

\begin{itemize}
  \item o \textbf{arredondamento} avalia a curvatura das arestas e dos vértices;
  \item a \textbf{esfericidade} avalia quanto a forma geral do grão se aproxima de uma esfera.
\end{itemize}

Um grão pode, portanto, ser bem arredondado e apresentar baixa esfericidade. A escala visual de Powers distingue seis classes de arredondamento: muito angular, angular, subangular, subarredondado, arredondado e bem arredondado \parencite{powers1953}.

\subsection{Arcabouço, matriz, cimento e poros}

Na descrição textural de uma rocha siliciclástica devem ser reconhecidos:

\begin{itemize}
  \item \textbf{arcabouço}: conjunto dos grãos que formam a estrutura principal;
  \item \textbf{matriz}: material detrítico mais fino depositado entre os grãos;
  \item \textbf{cimento}: mineral precipitado durante a diagênese;
  \item \textbf{poros}: espaços vazios que podem armazenar fluidos.
\end{itemize}

Matriz e cimento não são sinônimos: a matriz é deposicional ou sinsedimentar, enquanto o cimento se forma por precipitação diagenética.

\subsection{Exemplo de descrição textural}

Uma amostra composta por aproximadamente 80\% de areia fina, 10\% de areia média e 10\% de grânulos apresenta predomínio de areia fina, mas contém uma fração grossa relevante. A presença simultânea dessas classes sugere seleção baixa a moderada. A descrição deve ser completada com arredondamento, esfericidade, tipo de suporte, proporção de matriz e composição mineralógica, por exemplo, predomínio de quartzo e feldspato.
